% sections/08_progressivity.tex -- Section 8, Benefit Levels versus Progressivity (W3b).
% Plan: BLUEPRINT 2.8 and registry 8.9; Overrides O1 (no visible placeholders), O5 (provenance comments).
% Provenance abbreviations in "% src:" comments: T:x = paper/tables/x; 12:NN = frot/sub2025/12_app_quant.tex:NN;
% 7:NN = frot/sub2025/7_policy_reforms.tex:NN; 5:NN = frot/sub2025/5_quantitative_model.tex:NN.
% Every cell quoted below is in tables/welfare_progressivity_mvpfv7.tex (abbreviated "T:prog" after the first mention).
\section{Benefit Levels versus Progressivity}\label{sec:progressivity}

The reforms of \autoref{sec:claiming} change the claiming-age factor $\psi(a^c)$ in \eqref{eq:adjust} and leave the benefit formula unchanged. This section does the reverse: it raises the bracket rates of the primary insurance amount \eqref{eq:pia}, one bracket at a time or all three together, and leaves the claiming-age schedule at its status-quo values. The accounting is that of \autoref{sec:accounting}: the program outlay $B$ is retirement benefits, the Direct column is on the base $B^{0}$, and the Net column is on baseline net spending $G^{0}$ (\autoref{tab:conventions}).

\subsection{Reform grid}\label{subsec:pgrid}

Four sequences of reforms are considered; \autoref{tab:menu} lists them. Reforms P1.1--P1.9 raise the first bracket rate from 0.90 by one percentage point per step, to 0.91 under P1.1 and 0.99 under P1.9, and leave the second and third rates at 0.32 and 0.15. % src: 12:475; T:welfare_progressivity_mvpfv7.tex:115-123
Reforms P2.0--P2.9 start from P1.9 and raise the second bracket rate from 0.32 by 0.32 percentage points per step, one percent of its initial value, to 0.352 under P2.9. % src: 12:475; T:prog:126-135
Reforms P3.0--P3.9 start from P2.9 and raise the third bracket rate from 0.15 by 0.15 percentage points per step, again one percent of its initial value, to 0.165 under P3.9. % src: 12:475; T:prog:137-146
Reforms L0--L9 raise all three rates proportionally by one percent per step, from (0.909, 0.3232, 0.1515) under L0 to (0.99, 0.352, 0.165) under L9. % src: 12:475; T:prog:149, :158
The P sequence and the L sequence therefore begin at the status quo and end at the same policy---L9 and P3.9 are identical---but follow different paths: the P sequence takes the first bracket to its final value before touching the second, and the second before the third, whereas the L sequence moves the three rates together. % src: T:prog:146 vs :158
The bend points, the cap on average indexed earnings, the claiming-age schedule, the earnings test, and the payroll tax are unchanged.
% TODO-AUTHORS: disability benefits (eq:di) and pension income (eq:pension) are functions of the PIA in the model; confirm whether the reformed bracket rates also enter db_t and pb_t in the reform code. If so, add one sentence here saying that these flows change with the formula and are counted in O.

By \eqref{eq:pia}, a higher first-bracket rate raises the benefit of everyone with positive average indexed earnings: by a fixed dollar amount for individuals whose earnings exceed the first bend point and by a smaller amount below it, so that as a share of the benefit the increase is largest at low lifetime earnings. A higher second- or third-bracket rate raises benefits only for individuals whose average earnings exceed the first or the second bend point. A proportional increase raises every benefit by the same percentage. The three sequences thus place the same kind of instrument---a higher replacement rate---at different points of the lifetime-earnings distribution.

% tables/new_progressivity_compressed.tex -- Table 8 (tab:progressivity), written by W3b.
% Ten rows COPIED VERBATIM (sed, no retyping) from tables/welfare_progressivity_mvpfv7.tex:
%   P1.1 :115 | P1.5 :119 | P1.9 :123 | P2.4 :130 | P2.9 :135 | P3.4 :141 | P3.9 :146 | L0 :149 | L4 :153 | L9 :158
% Header rows copied from :112-114 (column layout unchanged: ID | bracket triple | MVPF | WTP | Net | Direct | Indir. | 30--49 | 62--64 | 65--69).
% Gradient macros (\gradXchange, \gradSSpay, \gradMVPF, \gradCost, \gradPP, \opacity) copied verbatim from :1-105 so that the
% shading is identical to Appendix Table C1 (the full file). Because every May 2025 table file defines these macros with
% \newcommand, they are undefined here before being (re)defined; the source table files are not edited.
% Caption and notes: BLUEPRINT 2.8; column definitions follow sections/05_accounting.tex and tables/new_accounting_conventions.tex.
\let\opacity\undefined \let\gradXchange\undefined \let\gradSSpay\undefined \let\gradMVPF\undefined
\let\gradCost\undefined \let\gradPP\undefined \let\gradMedpay\undefined \let\gradTaxchange\undefined
\begin{table}[!tp]
\centering
\caption{Social Security Progressivity and Level Reforms}
\label{tab:progressivity}
\vspace{6pt}
\setstretch{1.0}
\footnotesize
\setlength{\tabcolsep}{5pt}
\begin{threeparttable}
% ---- gradient macro block: verbatim copy of tables/welfare_progressivity_mvpfv7.tex:1-105 ----
% Color definitions
\definecolor{blue}{RGB}{0,114,178}
\definecolor{red}{RGB}{213,94,0}
\definecolor{white}{RGB}{255,255,255}
\newcommand*{\opacity}{100}% Full opacity

% Specific gradient commands for different columns
% For xchange% column: 0 is white, positive blue, negative red
\newcommand{\gradXchange}[1]{%
    \ifdim#1pt=0pt
        \cellcolor{white!\opacity} #1
    \else
        \ifdim#1pt>0pt
            % Positive values - blue
            \pgfmathparse{int(round(min(100, 100*(#1/1.5))))}% Scaling factor 1.5 based on max value in this column
            \xdef\tempa{\pgfmathresult}
            \cellcolor{blue!\tempa!white!\opacity} #1
        \else
            % Negative values - red (not present in this data but included for completeness)
            \pgfmathparse{int(round(min(100, 100*(-#1/1.5))))}% Same scaling as positive
            \xdef\tempa{\pgfmathresult}
            \cellcolor{red!\tempa!white!\opacity} #1
        \fi
    \fi
}

% For sspay %Δ column: 0 is white, positive red, negative blue
\newcommand{\gradSSpay}[1]{%
    \ifdim#1pt=0pt
        \cellcolor{white!\opacity} #1
    \else
        \ifdim#1pt>0pt
            % Positive values - red
            \pgfmathparse{int(round(min(100, 100*(#1/10.71))))}% Scaling based on max value ~10.71
            \xdef\tempa{\pgfmathresult}
            \cellcolor{red!\tempa!white!\opacity} #1
        \else
            % Negative values - blue (not present in this data but included for completeness)
            \pgfmathparse{int(round(min(100, 100*(-#1/10.71))))}
            \xdef\tempa{\pgfmathresult}
            \cellcolor{blue!\tempa!white!\opacity} #1
        \fi
    \fi
}

% For MVPF column: middle value is white (0.42), high values blue, low values red
\newcommand{\gradMVPF}[1]{%
    % Middle value is approximately 0.42 (2.82406434 * 0.15)
    \pgfmathparse{#1 - 0.42}
    \pgfmathsetmacro{\diff}{\pgfmathresult}
    
    \ifdim\diff pt=0pt
        \cellcolor{white!\opacity} #1
    \else
        \ifdim\diff pt>0pt
            % Higher than middle - blue
            \pgfmathparse{int(round(min(100, 100*(\diff/0.8))))}% Scaling factor based on range
            \xdef\tempa{\pgfmathresult}
            \cellcolor{blue!\tempa!white!\opacity} #1
        \else
            % Lower than middle - red
            \pgfmathparse{int(round(min(100, 100*(-\diff/1.8))))}% Scaling factor based on range
            \xdef\tempa{\pgfmathresult}
            \cellcolor{red!\tempa!white!\opacity} #1
        \fi
    \fi
}

% For total cost and indirect cost columns: 0 is white, positive red, negative blue
\newcommand{\gradCost}[1]{%
    \ifdim#1pt=0pt
        \cellcolor{white!\opacity} #1
    \else
        \ifdim#1pt>0pt
            % Positive values - red
            \pgfmathparse{int(round(min(100, 100*(#1/1.3))))}% Scaling based on max value
            \xdef\tempa{\pgfmathresult}
            \cellcolor{red!\tempa!white!\opacity} #1
        \else
            % Negative values - blue
            \pgfmathparse{int(round(min(100, 100*(-#1/0.37))))}% Scaling based on min value
            \xdef\tempa{\pgfmathresult}
            \cellcolor{blue!\tempa!white!\opacity} #1
        \fi
    \fi
}

% UPDATED: For age range p.p. columns: 0 is white, positive blue, negative red with more adapted scaling
\newcommand{\gradPP}[1]{%
    \ifdim#1pt=0pt
        \cellcolor{white!\opacity} #1
    \else
        \ifdim#1pt>0pt
            % Positive values - blue - adjusted scaling to better represent the positive range (max ~0.34)
            \pgfmathparse{int(round(min(100, 100*(#1/2.41))))}% Adjusted scaling for positive values
            \xdef\tempa{\pgfmathresult}
            \cellcolor{blue!\tempa!white!\opacity} #1
        \else
            % Negative values - red - adjusted scaling to better represent the negative range (min ~-2.41)
            \pgfmathparse{int(round(min(100, 100*(-#1/2.5))))}% Adjusted scaling for negative values
            \xdef\tempa{\pgfmathresult}
            \cellcolor{red!\tempa!white!\opacity} #1
        \fi
    \fi
}
% ---- end of macro block ----
% ---- header: verbatim copy of :108, :110-114 ----
\begin{tabular}{llllllllll}%
\toprule
    \hline
 \multicolumn{2}{c}{Policy Reforms} &\multicolumn{2}{c}{$\%\Delta$ Welfare}&\multicolumn{3}{c}{\%$\Delta$ Budget}&\multicolumn{3}{c}{$\Delta$ Participation}\\
        \hline
&PIA Brackets&MVPF&WTP&Net&Direct&Indir.&30--49&62--64&65--69\\	\hline
% ---- rows: verbatim copies of :115, :119, :123 ----
(P1.1)&(0.91,	0.32,	0.15)&\cellcolor{blue!90!white!\opacity}$[0.52,+\infty)$&\gradXchange{0.10}&\gradCost{-0.03}&\gradSSpay{0.57}&\gradCost{-0.10}&\gradPP{0.08}&\gradPP{0.20}&\gradPP{0.06}\\
(P1.5)&(0.95,	0.32,	0.15)&\gradMVPF{0.91}&\gradXchange{0.50}&\gradCost{0.21}&\gradSSpay{2.32}&\gradCost{-0.09}&\gradPP{0.16}&\gradPP{-0.34}&\gradPP{-0.31}\\
(P1.9)&(0.99,	0.32,	0.15)&\gradMVPF{0.69}&\gradXchange{0.80}&\gradCost{0.44}&\gradSSpay{4.12}&\gradCost{-0.11}&\gradPP{0.25}&\gradPP{-0.85}&\gradPP{-0.60}\\
\hline
% ---- rows: verbatim copies of :130, :135 ----
(P2.4)&(0.99,	0.336,	0.15)&\gradMVPF{0.38}&\gradXchange{0.92}&\gradCost{0.93}&\gradSSpay{7.21}&\gradCost{-0.05}&\gradPP{0.24}&\gradPP{-1.60}&\gradPP{-1.63}\\
(P2.9)&(0.99,	0.352,	0.15)&\gradMVPF{0.29}&\gradXchange{0.95}&\gradCost{1.27}&\gradSSpay{9.97}&\gradCost{-0.09}&\gradPP{0.28}&\gradPP{-2.26}&\gradPP{-2.05}\\
\hline
% ---- rows: verbatim copies of :141, :146 ----
(P3.4)&(0.99,	0.352,	0.1575)&\gradMVPF{0.29}&\gradXchange{0.96}&\gradCost{1.27}&\gradSSpay{10.34}&\gradCost{-0.13}&\gradPP{0.28}&\gradPP{-2.23}&\gradPP{-2.03}\\
(P3.9)&(0.99,	0.352,	0.165)&\gradMVPF{0.30}&\gradXchange{0.98}&\gradCost{1.24}&\gradSSpay{10.70}&\gradCost{-0.20}&\gradPP{0.28}&\gradPP{-2.37}&\gradPP{-1.99}\\
\hline \hline
% ---- rows: verbatim copies of :149, :153, :158 ----
(L0)&(0.909,	0.3232,	0.1515)&\gradMVPF{0.51}&\gradXchange{0.16}&\gradCost{0.12}&\gradSSpay{1.34}&\gradCost{-0.05}&\gradPP{0.11}&\gradPP{-0.05}&\gradPP{0.01}\\
(L4)&(0.945,	0.336,	0.1575)&\gradMVPF{0.59}&\gradXchange{0.74}&\gradCost{0.49}&\gradSSpay{5.28}&\gradCost{-0.21}&\gradPP{0.21}&\gradPP{-0.48}&\gradPP{-0.43}\\
(L9)&(0.99,	0.352,	0.165)&\gradMVPF{0.30}&\gradXchange{0.98}&\gradCost{1.24}&\gradSSpay{10.70}&\gradCost{-0.20}&\gradPP{0.28}&\gradPP{-2.37}&\gradPP{-1.99}\\
 \hline
 \bottomrule
    \end{tabular}
\begin{tablenotes}[flushleft]\scriptsize
\item[] \textit{Notes:} Ten of the 39 reforms of the grid in \autoref{subsec:pgrid}; the full grid is \autoref{tab:progressivity_full}. The second column gives the three bracket rates of the primary insurance amount \eqref{eq:pia} under the reform; the status quo is (0.9, 0.32, 0.15). L9 and P3.9 are the same policy.
% src: tables/welfare_progressivity_mvpfv7.tex:115-158 (39 rows); :146 vs :158 (L9 = P3.9); 5_quantitative_model.tex:102-104 (0.9/0.32/0.15)
Columns are defined in \autoref{tab:conventions} and \eqref{eq:ce}--\eqref{eq:mvpflb}: WTP is the consumption-equivalent welfare change in percent of consumption; Direct is the change in the net present value (NPV) of retirement benefits in percent of its baseline value; Indir.\ is the change in the NPV of all other taxes and transfers (disability benefits and transfers net of income and payroll taxes) in percent of its baseline value; Net is the change in net government spending in percent of its baseline value, related to the other two by \eqref{eq:netidentity} and not by addition; MVPF is the dollar welfare gain per dollar of net spending, and the entry $[0.52,+\infty)$ reports, for the one row in which net spending falls, the lower bound that counts only the increase in retirement benefits. The last three columns are changes in labor force participation at ages 30--49, 62--64, and 65--69 in percentage points. All NPVs are per capita at age 62 at the net return of 3.2 percent.
% src: 12_app_quant.tex:611 (SS definitions), :631-634 (MVPF, lower bound); 7_policy_reforms.tex:73 (participation groups); 5_quantitative_model.tex:390 (3.2 percent)
Shading is blue for a welfare gain, a fiscal saving, a higher MVPF, or a participation increase and red for the opposite, with intensity proportional to the entry. Source: simulations of the estimated model.
\end{tablenotes}
\end{threeparttable}
\end{table}


\subsection{Welfare, fiscal cost, and MVPF by bracket}\label{subsec:pmvpf}

\autoref{tab:progressivity} reports ten of the 39 rows of the grid: the first, a middle, and the last step of the first-bracket and level sequences, and a middle and the last step of the second- and third-bracket sequences; \autoref{tab:progressivity_full} reports all of them. Every reform in the grid raises welfare. The WTP rises from 0.10 percent of consumption under P1.1 to 0.80 under P1.9 and 0.98 under P3.9, and along the level sequence from 0.16 under L0 to 1.02 under L8 and 0.98 under L9. % src: T:prog:115, :123, :146, :149, :157, :158
The direct cost rises with generosity: the net present value of retirement benefits increases by 0.57 percent of its baseline value under P1.1, 4.12 under P1.9, 9.97 under P2.9, and 10.70 under P3.9 and L9. % src: T:prog:115, :123, :135, :146, :158
The indirect effect is a saving in every row: all other taxes and transfers fall by between 0.05 and 0.23 percent of their baseline value. % src: T:prog:115-158 (all Indir. entries negative; -0.05 at :127, :130, :149; -0.23 at :156)
Net government spending rises in every row except P1.1, for which it falls by 0.03 percent of its baseline value; P1.1 is approximately self-financing, and its MVPF is reported as the interval $[0.52,+\infty)$ of \eqref{eq:mvpflb}, where 0.52 is the welfare gain per dollar of additional retirement benefits. % src: T:prog:115

For the remaining rows the MVPF is finite and falls as the increase moves up the earnings distribution: it lies between 0.69 and 1.14 in the first bracket (P1.2--P1.9), falls from 0.56 under P2.0 to 0.29 under P2.9 in the second, and is 0.29 or 0.30 in every third-bracket row; along the level sequence it lies between 0.30 and 0.61. % src: T:prog:116-123 (minimum 0.69 at :123, maximum 1.14 at :120); :126, :135; :137-146; :149-158 (minimum 0.30 at :158, maximum 0.61 at :152)
The MVPF is not monotone within a sequence---P1.6 exceeds both P1.5 and P1.7, and L3 exceeds L2---but the ordering across sequences holds at every level of direct cost: a first-bracket row has a higher MVPF than the level row of comparable cost, and a second- or third-bracket row has no higher one. % src: T:prog:119-121 (0.91, 1.14, 1.01); :151-152 (0.49, 0.61); P3.7-P3.9 0.30 = L9 0.30 (:144-146, :158)
P1.9 and L3 raise outlays by 4.12 and 4.25 percent and have MVPFs of 0.69 and 0.61; P2.4 and L6 raise them by 7.21 and 7.41 percent and have MVPFs of 0.38 and 0.47; P2.9 and L8 raise them by 9.97 and 9.80 percent and have MVPFs of 0.29 and 0.35. % src: T:prog:123, :152, :130, :155, :135, :157

Within this menu, raising the first bracket rate is the most efficient way to increase the generosity of retirement benefits. Between P1.9 and P3.9 the WTP rises from 0.80 to 0.98 percent of consumption while the direct cost rises from 4.12 to 10.70 percent of outlays, so the spending added in the second and third brackets buys little additional welfare. % src: T:prog:123, :146
A uniform increase delivers less welfare per dollar of net spending than the same spending concentrated in the first bracket and more than the same spending placed in the upper brackets. The comparison is between increases in generosity under the same accounting, not between progressive and proportional benefit formulas in general. The pattern follows from where the transfer accrues. By \eqref{eq:pia} a first-bracket increase is, as a share of the benefit, largest for individuals with low lifetime earnings, who in the model also hold few assets in retirement and whose marginal utility of a transferred dollar is high in the sense of \autoref{sec:suffstats}; second- and third-bracket increases accrue only to individuals above the bend points, whose benefits are already larger. Changes in the level and progressivity of the benefit formula have been studied in general-equilibrium overlapping-generations models by \citet{nishiyama2007does}, \citet{Imrohoroglu2012}, and \citet{kitao2014sustainable}, in the life-cycle model of \citet{huggett2010well}, and in a life-cycle model with health and claiming by \citet{Yu2024}; the rows here are single-cohort, fixed-price experiments evaluated with the accounting of \autoref{sec:accounting}, and the comparison with those studies is qualitative.

\subsection{Participation responses}\label{subsec:ppart}

Participation at ages 30--49 rises in every row of the grid, from 0.08 percentage points under P1.1 to 0.28 under P2.9, under every third-bracket reform, and under L9. % src: T:prog:115 (0.08, the minimum), :135, :137-146, :158 (0.28, the maximum; also reached at :132 and :134)
Participation at ages 62--64 rises under P1.1 and P1.3, by 0.20 and 0.15 \pp, and falls in every other row, by as much as 2.41 \pp\ under P3.5. % src: T:prog:115, :117, :142
Participation at ages 65--69 is unchanged or higher under P1.1, P1.2, P1.3, and L0 (0.06, 0.00, 0.01, and 0.01 \pp) and lower in every other row, by as much as 2.1 \pp\ under P2.8. % src: T:prog:115-117, :149, :134
P1.1 and P1.3 are the only reforms that raise participation in all three age groups; from P1.4 onward in the P sequence and from L1 onward in the L sequence, participation falls at both older age groups in every row. % src: T:prog:115, :117 (all three positive); :118-146, :150-158 (62--64 and 65--69 negative)

The two margins respond to different components of the reform. At working ages a higher bracket rate raises, for individuals whose average indexed earnings fall in the affected bracket, the increment to future benefits from an additional year of earnings through \eqref{eq:aime} and \eqref{eq:pia}, a substitution effect toward work. At the claiming ages a higher benefit at every claiming age is an income effect toward earlier exit, and it grows with the size of the increase. The indirect column is the fiscal counterpart of these responses together with the induced changes in income taxes, payroll taxes, transfers, and disability benefits; it is a saving in every row of the grid, and its components are not reported. % src: T:prog:115-158 (Indir. negative in every row)
% TODO-AUTHORS: Open question 8 -- a decomposition of O (income taxes, payroll taxes, floor transfers, disability benefits) from existing output would support a mechanism sentence for the negative indirect effect. The May 2025 text described the opposite sign and a share attributed to income taxes from additional participants that has no cell; neither is reproduced. Without the decomposition only the printed sign is described.
The participation columns also qualify the MVPF ordering of \autoref{subsec:pmvpf}: the reforms with the highest MVPFs are those under which old-age participation falls least, and the upper-bracket and larger level reforms combine the smallest welfare gain per dollar with the largest declines in participation at ages 62--69. % src: T:prog:115-123 vs :126-158
